Unfortunately, estimates on boundary terms are rather delicate, so we will put off this subject until the next section. In the meantime, we will concentrate on interior estimates of higher-order derivatives. For example, let $\Omega^{\prime}$ be any domain such that $\Omega^{\prime}\subset\subset\Omega$. (The notation $\Omega^{\prime}\subset\subset\Omega$ means that $\Omega^{\prime}$ is compactly contained in the open set $\Omega$; i.e., $\overline{\Omega^{\prime}}$ is compact and $\overline{\Omega^{\prime}}\subset\Omega$.) We now choose a cutoff function $\zeta\in\mathcal{D}(\Omega)$ such that $0\leq\zeta\leq1$ and $\zeta\equiv1$ on $\Omega^{\prime}$. We can now make some calculations very similar to those above, but without any boundary terms getting in the way.

\[\begin{aligned}
\int_\Omega\zeta^2|\Delta u|^2\,\mathrm{d}\bm{x}
&=\sum_{i,j=1}^n\int_\Omega u_{x_ix_i}u_{x_jx_j}\zeta^2\,\mathrm{d}\bm{x}\\
&=-\sum_{i,j=1}^n\int_\Omega u_{x_i}u_{x_jx_ix_j}\zeta^2\,\mathrm{d}\bm{x}\\
&\quad-\sum_{i,j=1}^n\int_\Omega u_{x_i}u_{x_jx_j}2\zeta\zeta_{x_i}\,\mathrm{d}\bm{x}\\
&=\sum_{i,j=1}^n\int_\Omega u_{x_ix_j}u_{x_jx_i}\zeta^2\,\mathrm{d}\bm{x}\\
&\quad+\sum_{i,j=1}^n\int_\Omega u_{x_i}u_{x_jx_i}2\zeta\zeta_{x_j}-u_{x_i}u_{x_jx_j}2\zeta\zeta_{x_i}\,\mathrm{d}\bm{x}.
\end{aligned}\]

We now use this with inequalities of the form
\begin{equation}\label{eq:interior-cutoff-young-product-bound}
|u_{x_i}u_{x_ix_j}2\zeta\zeta_{x_j}|\leq\epsilon u_{x_ix_j}^2\zeta^2+\frac{1}{\epsilon}u_{x_i}^2\zeta_{x_j}^2
\end{equation}
and
\begin{equation}\label{eq:interior-cutoff-weighted-data-bound}
\int_\Omega|\Delta u|^2\zeta^2\,\mathrm{d}\bm{x}=\int_\Omega|f|^2\zeta^2\,\mathrm{d}\bm{x}\leq\int_\Omega|f|^2\,\mathrm{d}\bm{x}
\end{equation}
to get
\begin{equation}\label{eq:interior-cutoff-preabsorption-hessian-bound}
\int_\Omega|\nabla^2u|^2\zeta^2\,\mathrm{d}\bm{x}\leq\int_\Omega[|f|^2+\epsilon|\nabla^2u|^2\zeta^2+C(\epsilon)|\nabla u|^2|\nabla\zeta|^2]\,\mathrm{d}\bm{x}.
\end{equation}
We now let $\epsilon=1/2$ and use the fact that $\zeta\equiv1$ on $\Omega^{\prime}$ to get
\begin{equation}\label{eq:interior-cutoff-local-hessian-apriori-estimate}
\int_{\Omega^{\prime}}|\nabla^2u|^2\,\mathrm{d}\bm{x}\leq\int_\Omega\zeta^2|\nabla^2u|^2\,\mathrm{d}\bm{x}\leq C\left(\|f\|_2^2+\|\nabla u\|_2^2\right).
\end{equation}


Thus, we have an estimate on the $H^2(\Omega^{\prime})$ norm of a solution $u$ for any $\Omega^{\prime}\subset\subset\Omega$ in terms of the $L^2(\Omega)$ of the data $f$ and the $H^1(\Omega)$ norm of $u$. Of course, one of the major objections to the calculations performed above is that we needed to make unwarranted assumptions about the smoothness of the solution $u$ in order to perform the integrations by parts involved. In the rigorous version of these calculations below, these operations are replaced by analogous techniques involving difference quotients. Because the technique of using difference quotients is so important in this section, we present the following brief discussion of this topic.
